% Capstone (мини-курсовая) spec template – copy to NN-slug/spec.tex (see materials/sem2/capstone/README.md).
% In a real capstone, the numbers of "Требования к реализации" come from the generated
% requirements.tex (\input{requirements}); the template has none and types them.
\documentclass{ca-capstone}
\capstonerefs{sagemath,hps}
\begin{document}
\capstonetitle{Решение задачи DLP алгоритмом Полига"--~Хеллмана}

\goal
Требуется реализовать метод решения задачи дискретного логарифмирования в
$\GF{p} \cong \Zn{p}$ алгоритмом Полига"--~Хеллмана.

\functionality
\begin{enumerate}
  \item Решение DLP, в которой порядок основания "--- простое число, методом Шенкса или
    перебором\footnote{Перебором имеет смысл решать относительно <<маленькие>> DLP.}.
  \item Решение DLP, в которой порядок основания "--- степень простого числа.
  \item Решение DLP с составным порядком основания произвольного вида.
\end{enumerate}

\requirements
\begin{enumerate}
  \item Корректность решения задачи DLP.
  \item Приемлемая скорость работы: примеры с числами размером 50--100 бит "--- десятки секунд.
\end{enumerate}

\details
Рекомендуется использовать язык Python и систему компьютерной алгебры SageMath~\cite{sagemath}.

\schedule
Работа выполняется индивидуально. Срок выполнения "--- месяц.

\remarks
\begin{enumerate}
  \item Рекомендуется сначала реализовать метод Шенкса.
\end{enumerate}

\references
\end{document}
